\documentclass[11pt,a4paper]{article}  % Déclare la classe du document.

\newcommand{\chaptermark}[2]{\og #2 \fg{} (#1)} % ligne uniquement pour éviter une erreur en cas d'article
\include{../1ere-preambule}
\usepackage{mathrsfs}

\rhead{Exercices d'entrainement - Dérivation}

\newcommand{\va}[1]{\lvert \ #1 \ \rvert}
\newcommand{\vaf}[1]{\left\lvert \ #1\  \right\rvert}

\begin{document}
	\graphicspath{{img_entrainement/}}
	
	\begin{center}
		
		
		\fbox{
			\begin{minipage}[]{15cm}
				\begin{center}
					\LARGE{\rule{0.5cm}{0pt}\rule[-0.25cm]{0pt}{0.9cm}
						Exercices d'entrainement\\
						Dérivation
					}
				\end{center}
			\end{minipage}
		}
	\end{center}

\section*{Nombre dérivée}

\exercice{} 

Soit la fonction $f: x \mapsto(4-x)^2$ définie sur \R. Montrer que $f$ est dérivable en 1 et déterminer $f^{\prime}(1)$.

\exercice{}

Soit la fonction $f: x \mapsto\dfrac{-7}{3-x}$ définie sur $]-\infty;3[$.\\
Montrer que $f$ est dérivable en 2 et déterminer $f^{\prime}(2)$.

\exercice{}

Soit $f$ la fonction définie sur \R par $f: x \mapsto x^2+1$.
\begin{enumMet}
	\item Montrer que $f$  est dérivable en 3 et déterminer $f^{\prime}(3)$.
	\item Montrer que $f$ est dérivable en $-2$ et déterminer $f^{\prime}(-2)$.
\end{enumMet}

\exercice{}

Soit $f$ la fonction définie sur $\R^*$ par $f: x \mapsto \dfrac{3}{x}$.
\begin{enumMet}
	\item Montrer que $f$  est dérivable en -3 et déterminer $f^{\prime}(-3)$.
	\item Montrer que $f$ est dérivable en 1 et déterminer $f^{\prime}(1)$.
\end{enumMet}

\section*{Equation réduite d'une tangente}

\exercice{}

Soit une fonction $f$ définie et dérivable sur \R telle que $f(2)=5$ et $f^{\prime}(2)=-1$.
Soit $\mathscr{C}_f$ sa courbe représentative. Déterminer l'équation réduite de la tangente à $\mathscr{C}_f$ au point d'abscisse 2 .

\exercice{}

Soit une fonction $g$ définie et dérivable sur \R telle que $g(4)=-1$ et $g^{\prime}(4)=2$. Soit $\mathscr{C}_g$ sa courbe représentative. Déterminer l'équation réduite de la tangente à $\mathscr{C}_g$ au point d'abscisse 4.

\exercice{}

Soit $g$ la fonction définie sur \R par $g(x)=\left(2 x^2-5 x+4\right)^{10}$. On admet que $g$ est dérivable en 1 , et que $g^{\prime}(1)=-10$. Déterminer l'équation réduite de la tangente à la courbe représentative de $g$ au point d'abscisse 1.

\exercice{}

Soit $f$ une fonction dérivable sur $\mathbb{R}$ et soit $\mathscr{C}_f$ sa courbe représentative dans un repère du plan. Le point $A(3 ;-1)$ appartient à la courbe $\mathscr{C}_f$.
Sachant que la tangente $T$ à la courbe $\mathscr{C}_f$ au point $A$ passe par le point $J(0 ; 1)$, déterminer $f^{\prime}(3)$, puis l'équation de $T$.

\section*{Détermination graphique du nombre dérivé}

\begin{multicols}{2}

\exercice{}

Sur le graphique suivant, on donne la courbe représentative $\mathscr{C}_g$ d'une fonction $g$ sur l'intervalle $[-5 ; 8]$. Les droites $D_1, D_2$ et $D_3$ sont respectivement tangentes à $\mathscr{C}_g$ aux points A d'abscisse -3, B d'abscisse 2, et C d'abscisse 6 .
\begin{center}
	\includegraphics[scale=0.6]{1.05.03-Img01}
\end{center}
Lire sur le graphique les valeurs de $g(-3), g(2), g(6)$ et $g^{\prime}(-3), g^{\prime}(2), g^{\prime}(6)$

\exercice{}

La courbe d'une fonction $g$ définie sur $[-3 ; 5]$ est représentée ci-contre. La tangente à cette courbe au point $A$ d'abscisse 3 passe par le point de coordonnées $(-3 ; 6)$. Que vaut $g(3)$ ? Que vaut $g^{\prime}(3)$ ?
\begin{center}
	\includegraphics[scale=0.8]{1.05.03-Img03}
\end{center}

\exercice{}

Soit $\mathscr{C}_f$ la courbe représentative d'une fonction $f$ définie sur $\mathbb{R}$. $T_A$ est la tangente à $\mathscr{f}_f$ au point d'abscisse -2 . On sait que $f$ est dérivable sur $\mathbb{R}$ et que $f^{\prime}(4)=-4$.
\begin{center}
	\includegraphics[scale=0.8]{1.05.03-Img04}
\end{center}
\begin{enumMet}
	\item Déterminer $f^{\prime}(-2)$ et donner l'équation de la tangente $T_{\mathrm{A}}$. 
	\item B est le point de la courbe d'abscisse 4. Déterminer l'équation de la tangente à $\mathscr{C}_f$ au point B.
\end{enumMet}

\exercice{}

Dans le repère orthonormé $(0 ; \vv{i}, \vv{j})$ ci-contre, la courbe rouge $\mathscr{C}_f$ représente une fonction $f$ définie et dérivable sur $\mathbb{R}$, les droites tracées en bleu représentent les tangentes à $\mathscr{C}_f$ respectivement au point $\mathrm{O}$, au point $B$ d'abscisse 1 et au point C d'abscisse 3.
\begin{center}
	\includegraphics[scale=0.8]{1.05.03-Img05}
\end{center}
\begin{enumMet}
	\item Déterminer graphiquement $f^{\prime}(0), f^{\prime}(1)$ et $f^{\prime}(3)$.
	\item Déterminer l'équation réduite de la tangente à $\mathrm{C}_f$ au point $\mathrm{C}$.
	\item La courbe $\mathscr{C}_f$ est la représentation graphique de la fonction $f: x \mapsto x^3-4 x^2+2 x$
	Retrouver par le calcul les résultats des questions 1. et 2.
\end{enumMet}

\end{multicols}

\section*{Fonctions dérivées}

\exercice{}

Pour chacune des fonctions suivantes déterminer sur quel ensemble elle est dérivable puis déterminer sa dérivée.
\begin{multicols}{2}
\begin{enumMet}
	\item $f: x \mapsto 9 x^4$ définie sur $\mathbb{R}$
	\item $h: x \mapsto \dfrac{3}{4} x-7$ définie sur $\mathbb{R}$
	\item  $i: x \mapsto \sqrt{x}(3 x-1)$ définie sur $[0 ;+\infty[$
	\item $j: x \mapsto \dfrac{x^3}{1-x}$ définie sur $\mathbb{R} \backslash\{1\}$
\end{enumMet}
\end{multicols}

\exercice{}

Déterminer sur quel ensemble est dérivable chacune des fonctions suivantes, puis déterminer sa dérivée.
\begin{multicols}{2}
	\begin{enumMet}
		\item $f: x \mapsto-1,2 x^4+7 x^3-x$ définie sur $\mathbb{R}$
		\item $g: x \mapsto \dfrac{3}{5} x^5-\dfrac{1}{9} x^3-\dfrac{11}{2}$ définie sur $\mathbb{R}$
		\item $h: x \mapsto x^3(11-6 x)$ définie sur $\mathbb{R}$
		\item $j: x \mapsto \dfrac{25}{-10 x+9}$ définie sur $\mathbb{R} \backslash\left\{\dfrac{9}{10}\right\}$
	\end{enumMet}
\end{multicols}

\exercice{}

Déterminer sur quel ensemble est dérivable chacune des fonctions suivantes, puis déterminer sa dérivée.
\begin{multicols}{2}
	\begin{enumMet}
		\item $f: x \mapsto \dfrac{1}{2} x^2+\dfrac{1}{x^4}$
		\item $g: x \mapsto \dfrac{1}{x}(9-6 x)$
		\item $h: x \mapsto \dfrac{2}{1-x}$
		\item $i: x \mapsto \dfrac{x^2-1}{4 x+7}$
	\end{enumMet}
\end{multicols}


\exercice{}

Soit $g$ la fonction définie sur $\mathbb{R}$ par $g: x \mapsto(-9 x+1)^5$ et $h$ la fonction définie sur $[0 ;+\infty[$ par $h: x \mapsto \sqrt{3 x-1}$.\\
Pour les 2 fonctions, déterminer sur quel ensemble elle est dérivable puis déterminer sa dérivée.

\exercice{}

Déterminer pour chacune des fonctions suivantes, I'ensemble I sur lequel elle est dérivable, puis sa fonction dérivée sur $\mathrm{I}$.
\begin{multicols}{2}
	\begin{enumMet}
		\item $f(x)=\dfrac{5}{2 x}+\dfrac{3}{4}-\dfrac{7 x^2}{4}$
		\item $g(x)=\dfrac{-4}{5 x}(x-11)$
		\item $h(x)=\dfrac{5 x^2-8 x+1}{21-7 x}$
		\item $j(x)=\dfrac{-5}{3 x^2+2}$
		\item $k(x)=\dfrac{9 x}{x^2-6 x+5}$
		\item $m(x)=\sqrt{10-x}$
	\end{enumMet}
\end{multicols}


\exercice{}

Déterminer pour chacune des fonctions suivantes, l'ensemble I sur lequel elle est dérivable, puis sa fonction dérivée sur I.
\begin{multicols}{2}
	\begin{enumMet}
		\item $m: x \mapsto-\dfrac{2}{9} x^3+4 x^2-5$
		\item $n: x \mapsto-\dfrac{7 x}{2}(8 x+1)$
		\item $j: x \mapsto \dfrac{1}{5 x^5}$
		\item $p: x \mapsto-\dfrac{\sqrt{4 x+3}}{2}$
	\end{enumMet}
\end{multicols}

\end{document}

